\section{Introduction}

Germany continues to face a substantial shortage of workers in Science, Technology, Engineering, and Mathematics (STEM). According to the MINT-Herbstreport 2025, the aggregated STEM labour shortage amounted to 148,500 workers in October 2025, including 40,800 in academic STEM expert occupations \cite{Anger2025MINT}.

Science, Technology, Engineering, and Mathematics (STEM) disciplines offer favorable study and career prospects, also for people with disabilities. These fields place a strong emphasis on technical competencies and are characterized by highly structured content, which can be advantageous for diverse learning needs. Moreover, assistive technologies such as screen readers, speech recognition and output systems, as well as adapted workplace design and flexible work organization enable many people with disabilities to participate effectively in STEM education and employment.

Nevertheless, students with disabilities continue to encounter substantial barriers during their studies. These barriers arise from both structural and organizational conditions within higher education. Technical obstacles include limited accessibility of learning materials, especially those that rely heavily on visual representations or the use of complex graphical software tools. In addition, learning environments that are overly complex, poorly structured, or highly nested can impose a significant cognitive load on students with impairments such as attention-deficit/hyperactivity disorder, autism spectrum conditions, depression, or chronic fatigue.

As a result, learning progress may be slowed for students with disabilities, despite adequate or even high professional qualifications. Importantly, however, many of these barriers can be mitigated through relatively small and targeted measures that substantially improve accessibility. The patterns presented in this paper therefore aim to address key problems and challenges relevant to inclusive learning in STEM education.

\section{The Pattern Collection}

The pattern collection "A Pattern Collection for Generating Accessible Teaching Materials for Blind and Visually Impaired Students in Computer Science and Electrical Engineering" presented at EuroPLoP 2025 comprises five patterns:

\patternOneInline provides a general approach for gaining accessibility for \ac{bvi} students. 
The two patterns \patternTwoInline and \patternThreeInline introduce specialized solutions for different types of graphically represented teaching content. 
\patternFiveInline can be applied to gain accessibility of different types of graphics. 
\patternFourInline shows how time series can be made audible and thus accessible to \ac{bvi} people.

\newpage
\subsection*{Overview of the Pattern Collection from 2025}
\vspace{-2em}
%%%%%%%%%%%%%%%%% 1 1%%%%%%%%%%%%%%%%%%%%%%%%%%%%%%%%%
\begin{table}[!ht]
    \centering
    \begin{tabularx}{\textwidth}{|p{0.12\textwidth}|X|}
    \hline
         Pattern 1 &  \textsc{Core of Content}  \\ \hline
         Problem & Graphics are usually not accessible for the visually impaired.  Text within graphics cannot be extracted by assistive tools.  \\ \hline
         Solution & Extract the core content of teaching materials and represent it in a multi-directionally transformable format.  \\ \hline
    \end{tabularx}
\end{table}
\vspace{-3em}
%%%%%%%%%%%%%%%%% 2 %%%%%%%%%%%%%%%%%%%%%%%%%%%%%%%%%
\begin{table}[!ht]
    \centering
\begin{tabularx}{\textwidth}{|p{0.12\textwidth}|X|}
    \hline
         Pattern 2 & \textsc{Accessibility of Diagrams}  \\ \hline
         Problem &  Graphical representations of diagrams are not accessible to \ac{bvi} learners. \\ \hline
         Solution &  Analyze which concepts are represented by the diagram and find a textual representation for them. 
Implicit information as expressed by the spatial arrangement of diagram elements can be made tangible with haptic interface devices.  \\ \hline
\end{tabularx}
\end{table}
\vspace{-3em}
%%%%%%%%%%%%%%%%% 3 %%%%%%%%%%%%%%%%%%%%%%%%%%%%%%%%%
\begin{table}[!ht]
    \centering
\begin{tabularx}{\textwidth}{|p{0.12\textwidth}|X|}
    \hline
         Pattern  3 & \textsc{Accessibility of Graphics}  \\ \hline
         Problem &  Non-textual content—such as images, maps, charts, and diagrams remains largely inaccessible to \ac{bvi} students. This inaccessibility presents a significant barrier to information acquisition and equal participation in educational and professional environments.  \\ \hline
         Solution &  Provide textual, tactile, or auditory access to information represented visually in graphics.  \\ \hline
\end{tabularx}
\end{table}
\vspace{-3em}
%%%%%%%%%%%%%%%%% 4 %%%%%%%%%%%%%%%%%%%%%%%%%%%%%%%%%
\begin{table}[!ht]
    \centering
\begin{tabularx}{\textwidth}{|p{0.12\textwidth}|X|}
    \hline
         Pattern 4 & \textsc{Sonification of Time Series}  \\ \hline
         Problem & Graphical representations of time series  are not accessible for \ac{bvi} students and people in general. A sequential output of the data on Braille lines or by means of text-to-speech is not well suited to represent and convey the time based aspect of the data.   \\ \hline
         Solution &  Provide sonifications of time series. \\ \hline
\end{tabularx}
\end{table}
\vspace{-3em}
%%%%%%%%%%%%%%%%% 5 %%%%%%%%%%%%%%%%%%%%%%%%%%%%%%%%%
\begin{table}[!ht]
    \centering
\begin{tabularx}{\textwidth}{|p{0.12\textwidth}|X|}
    \hline
         Pattern 5 & \textsc{Graphics with Generated Explanation}  \\ \hline
         Problem &  Although electronic media can be made accessible to \ac{bvi} learners through the use of assistive technologies, this does not usually work for the graphic components embedded in these media. \\ \hline
         Solution &  Augment visually represented information with a generated and verified textual description.  \\ \hline
\end{tabularx}
\end{table}

This collection is extended by two additional patterns, \patternSixInline and \patternSevenInline. Figure~\ref{fig:pattern-relationships} shows Pattern~1 as the foundation of the collection. Patterns~2 to~5 provide specialized transformations of visual \ac{stem} artifacts. Pattern~6 uses and orchestrates these transformations in interactive learning situations, while Pattern~7 provides an audio-oriented overview and points learners to precise accessible representations where needed.

% KI: Robuste tabularx-Figur statt TikZ: kompiliert mit vorhandenen Paketen und bleibt vollständig editierbar.
\begin{figure}[ht]
\centering
\small
\setlength{\tabcolsep}{4pt}
\renewcommand{\arraystretch}{1.25}
\begin{tabularx}{\textwidth}{|p{0.18\textwidth}|X|}
\hline
\textbf{Layer 1: Foundation} & \textbf{Pattern 1: \textsc{Core of Content}} \\
\hline
\multicolumn{2}{|c|}{$\Downarrow$ foundation for all following patterns} \\
\hline
\textbf{Layer 2: Accessible representations} & \textbf{Patterns 2--5:} \textsc{Accessibility of Diagrams}; \textsc{Accessibility of Graphics}; \textsc{Sonification of Time Series}; \textsc{Graphics with Generated Explanation} \\
\hline
\textbf{Layer 3: AI-mediated support} & \textbf{Pattern 6: \textsc{Accessibility-Aware AI Tutor}} uses, orchestrates and explains Patterns 2--5 in interactive learning situations. \newline
\textbf{Pattern 7: \textsc{Structured Audio Companion}} summarizes and orients learners, and points to precise accessible representations where visual artifacts contain essential technical information. \\
\hline
\textbf{Verification relation} & Pattern 5 and Pattern 6 require verification and human-in-the-loop quality assurance when generated explanations are used. \\
\hline
\end{tabularx}
\caption{Pattern relationships in the extended collection. Pattern~1 provides the foundation. Patterns~2--5 provide specialized accessible representations for visual \ac{stem} artifacts. Pattern~6 uses these patterns for interactive AI-supported learning, while Pattern~7 offers audio-oriented orientation and points back to precise accessible representations.}
\label{fig:pattern-relationships}
\end{figure}
